%% 本文件由 工程文件/scripts/assemble.py 生成，勿手改。

%% begin units/ch10_u01.tex
% ch10_u01 —— 原书 p.375--p.380（第 10 章全部：10.1.1 -- 10.1.7）
% 原书章标题「10. To be inserted into Section 6, too」：本章是待插入第 6 节的材料（续）。

\bdchap{ch:ins6b}{待插入第 6 节的材料（续）}

\SourcePage{375}

\bdsec{10.1. Delta 函数}

本节的标题这样写可以吗？
% [Tutor Note: 原书 p.375 在节标题「10.1. Delta-functions.」之后紧接着印有作者手稿中的自问句「Is the title of the section OK ???」，底本如此，保留原样。]

\bdp{10.1.1} 由 8.1.5，有典范嵌入
\[
\mathfrak{d}_{\mathcal{L},\chi}\hookrightarrow\Gamma(\mathcal{GR},I_{\chi}\lambda_{\mathcal{L}}^{-1}).
\]
其像含于 $\Gamma(\mathcal{GR},I_{\chi}\lambda_{\mathcal{L}}^{-1})^{G(O)}$ 之中。李代数胚 $I/I^{2}$ 作用在 $\Gamma(\mathcal{GR},I_{\chi}\lambda_{\mathcal{L}}^{-1})^{G(O)}$ 上（见 ??? 与 5.6.8）。
% [Tutor Note: 原书 p.375 此处交叉引用印作「???」，是 LaTeX 未解析的引用；底本如此，保留原样。]
用 (81) 把 $I/I^{2}$ 与 3.5.11 中的李代数胚 $\mathfrak{a}_{{}^{L}\mathfrak{g}}$ 等同，其中 ${}^{L}\mathfrak{g}:=\operatorname{Lie}{}^{L}G$，而 ${}^{L}G$ 按 5.3.22 的意义理解。特别地，${}^{L}\mathfrak{g}$ 有一个指定的\footnote{§3 中我们用的是 $G$-oper 而非 ${}^{L}G$-oper，当时假定固定了一个 Borel 子群 $B\subset G$（见 3.1.1），故此处乐于有一个指定的 ${}^{L}\mathfrak{b}\subset{}^{L}\mathfrak{g}$。但事实上这并不本质：可以改写 §3 而不固定 $B$；此时只有李代数胚 $\mathfrak{a}_{\mathfrak{b}}$ 与 $\mathfrak{a}_{\mathfrak{n}}$，而没有具体的 $\mathfrak{b},\mathfrak{n}\subset\mathfrak{g}$。} Borel 子代数 ${}^{L}\mathfrak{b}$，以及一个指定的嘉当子代数 ${}^{L}\mathfrak{h}\subset{}^{L}\mathfrak{b}$；置 ${}^{L}\mathfrak{n}:=[{}^{L}\mathfrak{b},{}^{L}\mathfrak{b}]$。
% [Tutor Note: 原书 p.375 此处「Lie algebroid」印作「algenroid」，系排印笔误；底本如此，译文照译「李代数胚」，字形错误保留在注中。]
由 3.5.16，有李子代数胚 $\mathfrak{a}_{{}^{L}\mathfrak{n}}\subset\mathfrak{a}_{{}^{L}\mathfrak{b}}\subset\mathfrak{a}_{{}^{L}\mathfrak{g}}$，以及 $A_{{}^{L}\mathfrak{g}}(O)$-模的典范同构
\[
\mathfrak{a}_{{}^{L}\mathfrak{b}}/\mathfrak{a}_{{}^{L}\mathfrak{n}}=A_{{}^{L}\mathfrak{g}}(O)\otimes{}^{L}\mathfrak{h}.
\]
特别地，${}^{L}\mathfrak{h}\subset\mathfrak{a}_{{}^{L}\mathfrak{b}}/\mathfrak{a}_{{}^{L}\mathfrak{n}}$。

\bdp{10.1.2} \emph{定理.}
\begin{enumerate}[label=\roman*),leftmargin=2.2em]
\item $\mathfrak{a}_{{}^{L}\mathfrak{n}}$ 零化 $\mathfrak{d}_{\mathcal{L},\chi}$，故 $a\delta$ 对 $a\in{}^{L}\mathfrak{h}$、$\delta\in\mathfrak{d}_{\mathcal{L},\chi}$ 有意义。
\item 对 $a\in{}^{L}\mathfrak{h}$、$\delta\in\mathfrak{d}_{\mathcal{L},\chi}$，有 $a\delta=\chi(a)\delta$。
\end{enumerate}

\emph{注.} 我们把 $\chi\in P_{+}({}^{L}G)$ 与 ${}^{L}\mathfrak{h}$ 上的线性泛函等同，故 $\chi(a)$ 有意义。

陈述 (i) 不难。事实上，$\operatorname{Der}O$ 作用在 $\Gamma(\mathcal{GR},I_{\chi}\lambda_{\mathcal{L}}^{-1})^{G(O)}$ 上（见 5.6.10），而 $\mathfrak{a}_{{}^{L}\mathfrak{g}}$ 在该空间上的作用，与 $\operatorname{Der}O$ 在 $\mathfrak{a}_{{}^{L}\mathfrak{g}}$ 及该空间上的作用相容（用 (81) 的 $\operatorname{Der}O$-等变性，以及 3.6.16 末尾的注）。\footnote{事实上，有更强的结论：$\operatorname{Der}O$ 在 $\Gamma(\mathcal{GR},I_{\chi}\lambda_{\mathcal{L}}^{-1})^{G(O)}$ 上的作用，等于由 3.5.11 中定义的态射 $\operatorname{Der}O\to\mathfrak{a}_{{}^{L}\mathfrak{g}}$ 与 $\mathfrak{a}_{{}^{L}\mathfrak{g}}$ 在该空间上的作用复合而得者（这由 3.6.17 推出）。}于是陈述 (i) 由定理 8.1.4、命题 8.1.5 与 (77) 推出。类似地，用 (78) 可证：对 $a\in{}^{L}\mathfrak{h}$ 有 $a\mathfrak{d}_{\mathcal{L},\chi}\subset\mathfrak{d}_{\mathcal{L},\chi}$，但这弱于 (ii)。(ii) 将在\SourcePage{376} 10.1.3 – 10.1.7 中证明。在下面的证明中固定\footnote{$Z\operatorname{tors}_{\theta}(O)$ 的所有对象彼此同构。}$\mathcal{L}\in Z\operatorname{tors}_{\theta}(O)$，并把 $\lambda_{\mathcal{L}}$ 记作 $\lambda$，把 $\mathfrak{d}_{\mathcal{L},\chi}$ 记作 $\mathfrak{d}_{\chi}$，等等。

\bdp{10.1.3} 由 3.6.11，可把 10.1.2(ii) 改写成
\begin{equation}\tag{394}
a\delta=-(d(a),B\chi)\cdot\delta\quad\text{对 } a\in I^{\le 0},\ \delta\in\mathfrak{d}_{\chi}
\end{equation}
其中 $d:I^{\le 0}\to\mathfrak{h}$ 是 (83) 中的映射，$\chi$ 被看作 $\mathfrak{h}$ 的元素（见 10.1.2 的注），而 $B:\mathfrak{h}\to\mathfrak{h}^{*}$ 对应于内积 (18)。

\emph{注.} 「临界」内积 (18) 之所以出现在 (394) 的右端，是因为左端的定义涉及映射 (291)，而后者依赖于 $\mathfrak{g}$ 上内积的选取（见 5.6.11）。

\bdp{10.1.4} (394) 的证明方法将在 10.1.5 中描述。先说明必须克服的困难。$I/I^{2}$ 在 $\Gamma(\mathcal{GR},I_{\chi}\lambda^{-1})^{G(O)}$ 上的作用，来自重整化泛包络代数 $U^{\natural}$ 在 $\Gamma(\mathcal{GR},I_{\chi}\lambda^{-1})$ 上的作用，而后者由形变临界水平定义（见 ???）。
% [Tutor Note: 原书 p.376 此处交叉引用印作「???」，底本如此，保留原样。]
于是朴素的想法是形变 $I_{\chi}$，即构造一族 $\lambda^{h}$-扭 $\mathcal{D}$-模 $M^{?}_{h}$（$h\in\mathbb{C}$），使 $M^{?}_{0}=I_{\chi}$。但这不可能（至少在整体上），因为 $\operatorname{Orb}_{\chi}$ 上作为可逆 $\mathcal{O}$-模的 $\lambda^{h}$-扭 $\mathcal{D}$-模，只对一组离散的 $h$ 值存在。因此必须修改朴素的想法（见 10.1.5 与 10.1.7）。

\bdp{10.1.5} 下面使用 7.11.11 中 $\mathcal{D}_{\lambda^{h}}$-模的概念（于是 $h\in\mathbb{C}[h]$ 是参数）。在 10.1.7 中，将在 $U_{\chi}$ 上构造一个 $\mathcal{D}_{\lambda^{h}}$-模 $M$，以及嵌入
\begin{equation}\tag{395}
\mathfrak{d}_{\chi}\hookrightarrow\Gamma(U_{\chi},M\lambda^{-1})
\end{equation}
使得
\begin{enumerate}[label=(\roman*),leftmargin=2.2em]
\item $M$ 是平坦 $\mathbb{C}[h]$-模\footnote{于是对每个 $a\in\mathbb{C}$，有 $\mathcal{D}_{\lambda^{a}}:=\mathcal{D}_{\lambda^{h}}/(h-a)$ 上的模 $M_{a}:=M/(h-a)M$；可以说 $M$ 是由诸 $M_{a}$（$a\in\mathbb{C}$）构成的平坦族。}；
\item 有 $\mathcal{D}$-模态射 $M_{0}:=M/hM\to I^{U}_{\chi}:=I_{\chi}|_{U_{\chi}}$，使得合成
\[
\mathfrak{d}_{\chi}\hookrightarrow\Gamma(U_{\chi},M\lambda^{-1})\to\Gamma(U_{\chi},M_{0}\lambda^{-1})\to\Gamma(U_{\chi},I^{U}_{\chi}\lambda^{-1})
\]
等于 (392)；
\SourcePage{377}
\item (395) 的像被 $\mathfrak{g}\otimes\mathfrak{m}$ 零化，其中 $\mathfrak{m}$ 是 $O$ 的极大理想；
\item 对 $c\in C$（$C:=U\mathfrak{g}$ 的中心）与 $\delta\in\mathfrak{d}_{\chi}$，有
\begin{equation}\tag{396}
c\delta_{h}=\varphi(c)\delta_{h}
\end{equation}
其中 $\delta_{h}\in\Gamma(U_{\chi},M\lambda^{-1})$ 是 $\delta$ 在 (395) 下的像，$\varphi:C\to\mathbb{C}[h]$ 是对应于最高权为 $-hB\chi$ 的 Verma 模的特征标，而 $B:\mathfrak{h}\to\mathfrak{h}^{*}$ 是内积 (18)。
\end{enumerate}

\emph{注.} 1) $M\lambda^{-1}$ 是 $\mathcal{D}_{\lambda^{h+1}}$-模。2) 当然，$\mathcal{D}_{\lambda^{h+1}}:=\mathcal{D}_{\lambda^{s}}\otimes_{\mathbb{C}[s]}\mathbb{C}[h]$，其中态射 $\mathbb{C}[s]\to\mathbb{C}[h]$ 由 $s\mapsto h+1$ 定义。完全类似地可以定义例如 $\mathcal{D}_{\lambda^{-h}}$（此记号将在 10.1.7 中使用）。
% [Tutor Note: 原书 p.377 此处「similarly」印作「simialrly」，系排印笔误；底本如此，保留原样。]

\bdp{10.1.6} 下面从 (i) – (iv) 推出 (394)。由 5.6.7 – 5.6.8，(394) 的左端等于 $a^{\flat}\delta$，其中 $a^{\flat}\in U^{\flat}_{1}$ 与 $a\in I^{\le 0}$ 在 $U^{\flat}_{1}/U^{\flat}_{0}$ 中有相同的像。为构造 $a^{\flat}$，可把 $a$ 提升为元素 $\tilde a\in A$，其中 $A:=$ $\widetilde{\mathfrak{g}\otimes K}$ 的完备化泛包络代数，并使 $\tilde a$ 属于 $A$ 中由 $\mathfrak{g}\otimes O$ 拓扑生成的理想。于是 $h^{-1}\tilde a$ 属于 5.6.1 中的代数 $A^{\natural}$，故可令 $a^{\flat}:=$ $h^{-1}\tilde a$ 在 $U^{\natural}$ 中的像。下面将证明：对 $\tilde a$ 的适当选取\footnote{$a^{\flat}\delta$ 不依赖 $\tilde a$ 的选取，而 $a^{\flat}\delta_{0}$ 依赖（因为 $\delta_{0}$ 被 $\mathfrak{g}\otimes\mathfrak{m}$ 零化，但不被 $\mathfrak{g}\otimes\mathcal{O}$ 零化）。}，有
\begin{equation}\tag{397}
a^{\flat}\delta_{0}=-(d(a),B\chi)\cdot\delta_{0}
\end{equation}
\SourcePage{378}其中 $\delta_{0}$ 是 $\delta_{h}$ 在 $\Gamma(U_{\chi},M_{0}\lambda^{-1})$ 中的像，$d$、$B$ 的含义同 (394)。由 10.1.5(ii)，等式 (397) 蕴含 (394)。

下面描述 $\tilde a$ 的选取。可把 $a\in I^{\le 0}$ 写成 $c+a'$，其中 $c\in C$，而 $a'$ 属于 $\overline{U}'$ 中由 $\mathfrak{g}\otimes\mathfrak{m}$ 拓扑生成的左理想（按 3.6.8 – 3.6.9，有 $c=\pi(a)$）。我们取 $\tilde a\in A$，使 $\tilde a\mapsto a$，且 $\tilde a-c$ 属于 $A$ 中由 $\mathfrak{g}\otimes\mathfrak{m}$ 拓扑生成的左理想。于是 (397) 成立。

事实上，$M\lambda^{-1}$ 是 $\mathcal{D}_{\lambda^{h+1}}$-模。因此由 ???，$A^{\natural}$ 作用在 $\Gamma(U_{\chi},M\lambda^{-1})$ 上（能否就直接写 $M\lambda^{-1}$？），使得 $h:=\mathbf{1}-1\in\widetilde{\mathfrak{g}\otimes K}\subset A^{\natural}$ 按乘以 $h$ 作用（这样写可以吗？）。
% [Tutor Note: 原书 p.378 此处交叉引用印作「???」，是 LaTeX 未解析的引用；同句括号内的两个自问句原文作「can we write simply $M\lambda^{-1}$ ???」与「is this expression OK ???」，底本如此，均保留原样。]
于是可把 (397) 改写成
\begin{equation}\tag{398}
h^{-1}\tilde a\cdot\delta_{h}\equiv-(d(a),B\chi)\cdot\delta_{h}\mod h.
\end{equation}
由 10.1.5(iii) 与 10.1.5(iv)，有 $\tilde a\delta_{h}=c\delta_{h}=\varphi(c)\delta_{h}$。另一方面，$\varphi(c)\in\mathbb{C}[h]$ 模 $h^{2}$ 同余于 $-(d(a),B\chi)h$（见 10.1.5 中 $\varphi$ 的定义，以及 3.6.10 中 $d$ 的定义）。于是得到 (398)。

\bdp{10.1.7} 下面构造 $\mathcal{D}_{\lambda^{h}}$-模 $M$ 与态射 (395)，使之满足 10.1.5(i) – 10.1.5(iv)。

我们有 $\mathcal{GR}$ 上的 $G(O)$-等变线丛 $\lambda=\lambda_{\mathcal{L}}$。记 $\lambda_{\chi}$ 为它在 $\operatorname{Orb}_{\chi}$ 上的限制。$\operatorname{orb}_{\chi}$ 与 $p:\operatorname{Orb}_{\chi}\to\operatorname{orb}_{\chi}$ 的含义同 9.1.3 与 9.1.5。回忆 $G(O)$ 通过 $G(O/tO)=G$ 作用在 $\operatorname{orb}_{\chi}$ 上，且 $p$ 是 $G(O)$-等变的。由 9.1.11，存在 $\operatorname{orb}_{\chi}$ 上唯一的 $G$-等变线丛 $\underline{\lambda}_{\chi}$，使得 $\lambda_{\chi}=p^{*}\underline{\lambda}_{\chi}$。

在 $\operatorname{orb}_{\chi}$ 上有扭微分算子层 $\mathcal{D}_{\underline{\lambda}_{\chi}^{h}}$。置 $N:=p^{\dagger}\mathcal{D}_{\underline{\lambda}_{\chi}^{-h}}$，其中 $\mathcal{D}_{\underline{\lambda}_{\chi}^{-h}}$ 被看作左 $\mathcal{D}_{\underline{\lambda}_{\chi}^{-h}}$-模，而 $p^{\dagger}$ 是通常的拉回函子。$N$ 是 $\operatorname{Orb}_{\chi}$ 上的左 $\mathcal{D}_{\underline{\lambda}_{\chi}^{-h}}$-模，带有典范截面 $\mathbb{1}:=p^{\dagger}(1)\in\Gamma(\operatorname{Orb}_{\chi},N)$。显然 $\omega_{\operatorname{Orb}_{\chi}}\otimes_{\mathcal{O}}N$ 是 $\operatorname{Orb}_{\chi}$ 上的右\SourcePage{379}$\mathcal{D}_{\underline{\lambda}_{\chi}^{h}}$-模\footnote{顺便指出，$\omega_{\operatorname{Orb}_{\chi}}\otimes_{\mathcal{O}}N$ 典范同构于右 $\mathcal{D}_{\underline{\lambda}_{\chi}^{h}}$-模 $\omega_{\operatorname{orb}_{\chi}}\otimes_{\mathcal{O}}\mathcal{D}_{\underline{\lambda}_{\chi}^{h}}$ 的拉回。事实上，$\omega_{\operatorname{orb}_{\chi}}\otimes_{\mathcal{O}}\mathcal{D}_{\underline{\lambda}_{\chi}^{h}}$ 在把右 $\mathcal{D}_{\underline{\lambda}_{\chi}^{h}}$-模变为左 $\mathcal{D}_{\underline{\lambda}_{\chi}^{-h}}$-模的通常函子 $M\mapsto M\otimes_{\mathcal{O}}\omega_{\operatorname{orb}_{\chi}}^{-1}$ 下的像，由 $1\in\Gamma(\operatorname{orb}_{\chi},\omega_{\operatorname{orb}_{\chi}}\otimes_{\mathcal{O}}\mathcal{D}_{\underline{\lambda}_{\chi}^{h}}\otimes_{\mathcal{O}}\omega_{\operatorname{orb}_{\chi}}^{-1})$ 自由生成，故典范同构于 $\mathcal{D}_{\underline{\lambda}_{\chi}^{-h}}$。}。该截面 $\mathbb{1}$ 诱导一个 $\mathcal{O}$-模态射
\begin{equation}\tag{399}
\omega_{\operatorname{Orb}_{\chi}}\to\omega_{\operatorname{Orb}_{\chi}}\otimes_{\mathcal{O}}N.
\end{equation}

我们把 $M$ 定义为 $\omega_{\operatorname{Orb}_{\chi}}\otimes_{\mathcal{O}}N$ 沿闭嵌入 $\operatorname{Orb}_{\chi}\hookrightarrow U_{\chi}$ 的正向像。态射 (395) 定义为如下合成：
\[
\mathfrak{d}_{\chi}\hookrightarrow\Gamma(\operatorname{Orb}_{\chi},\omega_{\operatorname{Orb}_{\chi}}\otimes\lambda_{\chi}^{-1})\hookrightarrow\Gamma(\operatorname{Orb}_{\chi},(\omega_{\operatorname{Orb}_{\chi}}\otimes_{\mathcal{O}}N)\lambda_{\chi}^{-1})\hookrightarrow\Gamma(U_{\chi},M\lambda^{-1}),
\]
其中第一个态射由 (241) 诱导，第二个由 (399) 诱导。

性质 10.1.5(i) 是显然的。性质 10.1.5(ii) 也是显然的：态射 $M_{0}\to I^{U}_{\chi}$ 来自 $\mathcal{D}$-模态射 $N_{0}=p^{\dagger}\mathcal{D}_{\operatorname{orb}_{\chi}}\to\mathcal{O}_{\operatorname{Orb}_{\chi}}$，后者使 $\mathbb{1}\mapsto 1$（是否可以只写 $\mathbb{1}$，而不写 $\mathbb{1}\bmod h$ 或 $\mathbb{1}_{0}$ 等等？）。注意到 10.1.5(iii) 与 10.1.5(iv) 是关于 $\mathfrak{g}\otimes O$ 在 (395) 的像上的作用的性质。该像含于 $\mathfrak{g}\otimes O$-不变子空间（还是 $\mathbb{C}[h]$-子模？）
\begin{equation}\tag{400}
\Gamma(\operatorname{Orb}_{\chi},(\omega_{\operatorname{Orb}_{\chi}}\otimes_{\mathcal{O}}N)\lambda_{\chi}^{-1})=\Gamma(\operatorname{Orb}_{\chi},\lambda_{\chi}^{-1}\omega_{\operatorname{Orb}_{\chi}}\otimes_{\mathcal{O}}N).
\end{equation}
故为证明 10.1.5(iii) 与 10.1.5(iv)，只需在 $\operatorname{Orb}_{\chi}$ 上工作。利用 (241)，把 (400) 与
\begin{equation}\tag{401}
\mathfrak{d}_{\chi}\otimes\Gamma(\operatorname{Orb}_{\chi},N)
\end{equation}
等同。(400) 与 (401) 之间的同构是 $\mathfrak{g}\otimes O$-等变的（$\mathfrak{g}\otimes O$ 在 $\mathfrak{d}_{\chi}$ 上的作用平凡），因为同构 (241) 是 $\mathfrak{g}\otimes O$-等变的。故 10.1.5(iii) 与 10.1.5(iv) 等价于 $\mathbb{1}\in\Gamma(\operatorname{Orb}_{\chi},N)$ 的如下性质：\SourcePage{380}
\begin{equation}\tag{402}
(\mathfrak{g}\otimes\mathfrak{m})\mathbb{1}=0,
\end{equation}
\begin{equation}\tag{403}
c\mathbb{1}=\varphi(c)\mathbb{1}\quad\text{对 } c\in C.
\end{equation}
% [Tutor Note: 原书 p.379 此段两处括号是作者手稿中的自问句，原文作「(is it OK to write $\mathbb{1}$ instead of $\mathbb{1}$ mod $h$, or $\mathbb{1}_0$, etc. ???)」与「(or $\mathbb{C}[h]$-submodule ???)」，底本如此，保留原样。]

回忆 $C:=$ $U\mathfrak{g}$ 的中心，$\varphi:C\to\mathbb{C}[h]$ 表示对应于最高权为 $-hB\chi$ 的 Verma 模的特征标，而 $B:\mathfrak{h}\to\mathfrak{h}^{*}$ 是内积 (18)。

于是只剩证明 (402) 与 (403)。回忆 $N:=p^{\dagger}\mathcal{D}_{\underline{\lambda}_{\chi}^{-h}}$，$\mathbb{1}:=p^{\dagger}(1)$，且 $p:\operatorname{Orb}_{\chi}\to\operatorname{orb}_{\chi}$ 是 $G(O)$-等变的。因此 (402) 是显然的（因为 $\mathfrak{g}\otimes\mathfrak{m}$ 在 $(\operatorname{orb}_{\chi},\underline{\lambda}_{\chi})$ 上的作用平凡），而 (403) 等价于图
\begin{equation}\tag{404}
\begin{tikzcd}[row sep=2.6em,column sep=3.6em]
C \arrow[r, hook] \arrow[d, "\varphi"'] & U\mathfrak{g} \arrow[d] \\
\mathbb{C}[h] \arrow[r, hook] & \Gamma(\operatorname{orb}_{\chi},\mathcal{D}_{\underline{\lambda}_{\chi}^{-h}})
\end{tikzcd}
\end{equation}
的交换性。

回忆 $\underline{\lambda}_{\chi}$ 是 $\operatorname{orb}_{\chi}$ 上的 $G$-等变线丛，使得 $\lambda_{\chi}=p^{*}\underline{\lambda}_{\chi}$。由于 $\operatorname{orb}_{\chi}=G/P_{\chi}^{-}$（见 9.1.3），$\underline{\lambda}_{\chi}$ 的同构类由某个 $l\in\operatorname{Hom}(P_{\chi}^{-},\mathbb{G}_{m})\subset\operatorname{Hom}(H,\mathbb{G}_{m})\subset\mathfrak{h}^{*}$ 确定。事实上，
\begin{equation}\tag{405}
l=B\chi.
\end{equation}

确实，存在 $G(O)$-等变同构 $\lambda_{\chi}=\omega_{\operatorname{Orb}_{\chi}}$（见 (241)），故 $\underline{\lambda}_{\chi}$ 与 9.1.12 中的线丛 $\mathcal{M}_{\chi}$ 是 $G$-同构的，而 (405) 等价于命题 9.1.13。(404) 的交换性由 (405) 推出（见 ???）。
% [Tutor Note: 原书 p.380 此处交叉引用印作「???」，是 LaTeX 未解析的引用；底本如此，保留原样。]
证毕。
%% end units/ch10_u01.tex
